% GENERATED FILE. Edit canonical lessons, then run npm run generate:books.
% canonical-source-hash: fnv1a64:1db1d6c96a2a363c
% canonical-lessons: GU-C433-posak, GU-C433-jholi, GU-C433-kothlo, GU-C433-sobhvu, GU-C433-sankocavu

\chapter{Outfits and Bags}
\label{ch:gu-outfits-and-bags}
\hlchaptermodality{433}

\begin{chapteropening}
\textbf{By the end of this chapter:} \emph{I can talk about what suits me and what shrank.}

Complete the last lesson of chapter 433: I can talk about what suits me and what shrank.
\end{chapteropening}

\section[\texorpdfstring{\gu{પોશાક} (pośāk)}{પોશાક (pośāk)}]{\gu{પોશાક} (pośāk) — an outfit; a costume}

\label{lesson:GU-C433-posak}



\begin{quote}
Before the new one: say the Gujarati for to rub, then the Gujarati for to make shine.
\end{quote}

\subsection*{You'll want to know: \gu{પોશાક}}
\textbf{\gu{પોશાક}} — \emph{pośāk} — \textquotedblleft{}an outfit; a costume\textquotedblright{}.

\begin{quote}
\gu{લગ્નમાં} \gu{તેણે} \gu{નવો} \gu{પોશાક} \gu{પહેર્યો}.
\end{quote}

\begin{quote}
\emph{lagnamā̃ teṇe navo pośāk paheryo.}
\end{quote}

\begin{quote}
He wore a new outfit at the wedding.
\end{quote}

A masculine noun.

\begin{cousinweb}
\textbf{\gu{પોશાક}} (\emph{pośāk}) is usually traced to Persian \emph{pōšāk}, \textquotedblleft{}clothing, dress\textquotedblright{}.
\end{cousinweb}

\begin{grammarlens}[title={What you've built}]
One more word for getting around.
\end{grammarlens}

\subsection*{Your turn}
\begin{itemize}
\raggedright
  \item \emph{Say it:} \emph{pośāk}
  \item \emph{Say it:} \emph{pośāk}, once more
  \item \emph{Say it:} \emph{lagnamā̃ teṇe navo pośāk paheryo.} — He wore a new outfit at the wedding.
  \item \emph{Recall:} say the Gujarati for to rub, then the Gujarati for to make shine, then say \emph{pośāk} again
\end{itemize}

\begin{tcolorbox}[breakable,colback=teal!4,colframe=teal!35!black,title={Before you move on}]
What does \gu{પોશાક} mean? (\textquotedblleft{}An outfit; a costume\textquotedblright{}.) Say it once more.
\end{tcolorbox}
\section[\texorpdfstring{\gu{ઝોળી} (jhoḷī)}{ઝોળી (jhoḷī)}]{\gu{ઝોળી} (jhoḷī) — a shoulder bag (of cloth)}

\label{lesson:GU-C433-jholi}



\begin{quote}
Before the new one: say the Gujarati for to make shine, then the Gujarati for an outfit.
\end{quote}

\subsection*{You'll want to know: \gu{ઝોળી}}
\textbf{\gu{ઝોળી}} — \emph{jhoḷī} — \textquotedblleft{}a shoulder bag (of cloth)\textquotedblright{}.

\begin{quote}
\gu{દાદાએ} \gu{ઝોળીમાં} \gu{ફળ} \gu{ભર્યાં}.
\end{quote}

\begin{quote}
\emph{dādāe jhoḷīmā̃ phaḷ bharyā̃.}
\end{quote}

\begin{quote}
Grandfather filled his shoulder bag with fruit.
\end{quote}

A feminine noun.

\begin{cousinweb}
\textbf{\gu{ઝોળી}} (\emph{jhoḷī}) goes back to \emph{*jhōla}, a form CDIAL reconstructs; CDIAL lists the Gujarati word, \textquotedblleft{}a bag\textquotedblright{} (CDIAL 5415).
\end{cousinweb}

\begin{grammarlens}[title={What you've built}]
One more word for getting around.
\end{grammarlens}

\subsection*{Your turn}
\begin{itemize}
\raggedright
  \item \emph{Say it:} \emph{jhoḷī}
  \item \emph{Say it:} \emph{jhoḷī}, once more
  \item \emph{Say it:} \emph{dādāe jhoḷīmā̃ phaḷ bharyā̃.} — Grandfather filled his shoulder bag with fruit.
  \item \emph{Recall:} say the Gujarati for to make shine, then the Gujarati for an outfit, then say \emph{jhoḷī} again
\end{itemize}

\begin{tcolorbox}[breakable,colback=teal!4,colframe=teal!35!black,title={Before you move on}]
What does \gu{ઝોળી} mean? (\textquotedblleft{}A shoulder bag (of cloth)\textquotedblright{}.) Say it once more.
\end{tcolorbox}
\section[\texorpdfstring{\gu{કોથળો} (kothḷo)}{કોથળો (kothḷo)}]{\gu{કોથળો} (kothḷo) — a sack}

\label{lesson:GU-C433-kothlo}



\begin{quote}
Before the new one: say the Gujarati for an outfit, then the Gujarati for a shoulder bag.
\end{quote}

\subsection*{You'll want to know: \gu{કોથળો}}
\textbf{\gu{કોથળો}} — \emph{kothḷo} — \textquotedblleft{}a sack\textquotedblright{}.

\begin{quote}
\gu{હમાલ} \gu{ઘઉંનો} \gu{ભારે} \gu{કોથળો} \gu{ઊંચકે} \gu{છે}.
\end{quote}

\begin{quote}
\emph{hamāl ghaũno bhāre kothḷo ū̃cke che.}
\end{quote}

\begin{quote}
The porter lifts a heavy sack of wheat.
\end{quote}

A masculine noun.

\begin{cousinweb}
\textbf{\gu{કોથળો}} (\emph{kothḷo}) goes back to \emph{*kōtthala}, a form CDIAL reconstructs; CDIAL lists the Gujarati word, \textquotedblleft{}a large bag\textquotedblright{} (CDIAL 3511).
\end{cousinweb}

\begin{grammarlens}[title={What you've built}]
One more word for getting around.
\end{grammarlens}

\subsection*{Your turn}
\begin{itemize}
\raggedright
  \item \emph{Say it:} \emph{kothḷo}
  \item \emph{Say it:} \emph{kothḷo}, once more
  \item \emph{Say it:} \emph{hamāl ghaũno bhāre kothḷo ū̃cke che.} — The porter lifts a heavy sack of wheat.
  \item \emph{Recall:} say the Gujarati for an outfit, then the Gujarati for a shoulder bag, then say \emph{kothḷo} again
\end{itemize}

\begin{tcolorbox}[breakable,colback=teal!4,colframe=teal!35!black,title={Before you move on}]
What does \gu{કોથળો} mean? (\textquotedblleft{}A sack\textquotedblright{}.) Say it once more.
\end{tcolorbox}
\section[\texorpdfstring{\gu{શોભવું} (śobhvũ)}{શોભવું (śobhvũ)}]{\gu{શોભવું} (śobhvũ) — to suit; to look good (on someone)}

\label{lesson:GU-C433-sobhvu}



\begin{quote}
Before the new one: say the Gujarati for a shoulder bag, then the Gujarati for a sack.
\end{quote}

\subsection*{You'll want to know: \gu{શોભવું}}
\textbf{\gu{શોભવું}} — \emph{śobhvũ} — \textquotedblleft{}to suit; to look good (on someone)\textquotedblright{}.

\begin{quote}
\gu{આ} \gu{લાલ} \gu{સાડી} \gu{તને} \gu{ઘણી} \gu{શોભે} \gu{છે}.
\end{quote}

\begin{quote}
\emph{ā lāl sāḍī tane ghaṇī śobhe che.}
\end{quote}

\begin{quote}
This red sari suits you very well.
\end{quote}

Past (it …): \gu{તે} \gu{સાડી} \gu{તેને} \gu{શોભી}, \emph{te sāḍī tene śobhī}.

\begin{cousinweb}
\textbf{\gu{શોભવું}} (\emph{śobhvũ}) is built on Sanskrit \emph{śobhā}, \textquotedblleft{}splendour, beauty, grace\textquotedblright{} (Monier-Williams): to look beautiful in something.
\end{cousinweb}

\begin{grammarlens}[title={What you've built}]
One more word for getting around.
\end{grammarlens}

\subsection*{Your turn}
\begin{itemize}
\raggedright
  \item \emph{Say it:} \emph{śobhvũ}
  \item \emph{Say it:} \emph{śobhvũ}, once more
  \item \emph{Say it:} \emph{ā lāl sāḍī tane ghaṇī śobhe che.} — This red sari suits you very well.
  \item \emph{Recall:} say the Gujarati for a shoulder bag, then the Gujarati for a sack, then say \emph{śobhvũ} again
\end{itemize}

\begin{tcolorbox}[breakable,colback=teal!4,colframe=teal!35!black,title={Before you move on}]
What does \gu{શોભવું} mean? (\textquotedblleft{}To suit; to look good (on someone)\textquotedblright{}.) Say it once more.
\end{tcolorbox}
\section[\texorpdfstring{\gu{સંકોચાવું} (sankocāvũ)}{સંકોચાવું (sankocāvũ)}]{\gu{સંકોચાવું} (sankocāvũ) — to shrink}

\label{lesson:GU-C433-sankocavu}



\begin{quote}
Before the new one: say the Gujarati for a sack, then the Gujarati for to suit.
\end{quote}

\subsection*{You'll want to know: \gu{સંકોચાવું}}
\textbf{\gu{સંકોચાવું}} — \emph{sankocāvũ} — \textquotedblleft{}to shrink\textquotedblright{}.

\begin{quote}
\gu{ધોયા} \gu{પછી} \gu{મારું} \gu{ખમીસ} \gu{સંકોચાઈ} \gu{ગયું}.
\end{quote}

\begin{quote}
\emph{dhoyā pachī mārũ khamīs sankocāī gayũ.}
\end{quote}

\begin{quote}
My shirt shrank after it was washed.
\end{quote}

Past (it …): \gu{તે} \gu{ખમીસ} \gu{સંકોચાઈ} \gu{ગયું}, \emph{te khamīs sankocāī gayũ}.

\begin{cousinweb}
\textbf{\gu{સંકોચાવું}} (\emph{sankocāvũ}) is built on Sanskrit \emph{saṁkoca}, \textquotedblleft{}contraction, shrinking together\textquotedblright{} (Monier-Williams).
\end{cousinweb}

\begin{grammarlens}[title={What you've built}]
One more word for getting around.
\end{grammarlens}

\subsection*{Your turn}
\begin{itemize}
\raggedright
  \item \emph{Say it:} \emph{sankocāvũ}
  \item \emph{Say it:} \emph{sankocāvũ}, once more
  \item \emph{Say it:} \emph{dhoyā pachī mārũ khamīs sankocāī gayũ.} — My shirt shrank after it was washed.
  \item \emph{Recall:} say the Gujarati for a sack, then the Gujarati for to suit, then say \emph{sankocāvũ} again
\end{itemize}

\begin{tcolorbox}[breakable,colback=teal!4,colframe=teal!35!black,title={Before you move on}]
What does \gu{સંકોચાવું} mean? (\textquotedblleft{}To shrink\textquotedblright{}.) Say it once more.
\end{tcolorbox}
